\documentclass[10pt,a4paper]{amsart}
\usepackage[margin=19mm]{geometry}
\usepackage{amsmath,amssymb,mathtools}
\usepackage[scr=rsfs,cal=euler]{mathalfa}
\usepackage{xfrac}
\usepackage[unicode,hidelinks,bookmarks=false]{hyperref}
\newcommand{\ie}{i.e.\ }
\newcommand{\cf}{cf.\ }
\newcommand{\resp}{respectively\ }
\newcommand{\Subsection}{\subsection}
\providecommand{\nobreakdash}{\nobreak\hbox{-}}
\setlength{\emergencystretch}{2em}
\multlinegap0pt
\usepackage{bm}
\usepackage{bbm}
\usepackage{dsfont}
\usepackage{xspace}
\usepackage{cancel}
\usepackage{mathtools}
\usepackage{amsthm}
\makeatletter
\@ifundefined{theorem}{\newtheorem{theorem}{Theorem}}{}
\@ifundefined{lemma}{\newtheorem{lemma}[theorem]{Lemma}}{}
\@ifundefined{proposition}{\newtheorem{proposition}[theorem]{Proposition}}{}
\makeatother
\title[Ryll-Wojtaszczyk Formulas for bihomogeneous polynomials on t]{Ryll-Wojtaszczyk Formulas for bihomogeneous polynomials on the sphere}
\author{Andreas Defant, Daniel Galicer, Mart\'in Mansilla, Mieczys{\l}aw Masty{\l}o, Santiago Muro}
\date{Source: arXiv:2411.12579v2 (CC BY 4.0)}
\begin{document}
\maketitle

\noindent\emph{Source and licence.} Ryll-Wojtaszczyk Formulas for bihomogeneous polynomials on the sphere. Version arXiv:2411.12579v2 (CC BY 4.0), distributed under the Creative Commons Attribution 4.0 International licence, as recorded in the arXiv licence metadata for this exact version and shown on the abstract page https://arxiv.org/abs/2411.12579v2. Full rights notice: licenses/arxiv-2411.12579-CC-BY-4.0.txt.\par\smallskip

\noindent\textit{Editorial scope.} Counted unit: the Theorem whose source label is \texttt{Koornwinder} in the article (source0.tex lines 1432--1441, source label \texttt{Koornwinder}), delivered together with the article's own complete original proof of it (source0.tex lines 1444--1535, one \texttt{proof} environment of 92 lines). No mathematics is written, repaired or re-derived: statement and proof are the article's own text line for line. Required context retained uncounted: lemma: invariance (lines 359--362); propleg-comp (lines 1196--1205); repro (lines 1280--1284); integration (lines 1426--1429). Every definition, display and label of those lines is retained unchanged. Omitted from this package and not redistributed: 1--358, 363--1195, 1206--1279, 1285--1425, 1430--1431, 1442--1443, 1536--2553. Nothing inside a retained excerpt is omitted or altered: every retained body is delivered exactly as the article's lines are written, so no omission receipt of any kind accompanies this package. Citations: the retained excerpts carry no citation command at all, so no bibliography entry of the article is transcribed and no reference is redistributed. Reference adaptation: none was needed and none was made; every reference inside the delivered unit resolves inside it, so no unresolved reference or citation is produced. Printed slips kept literal: no printed word, symbol, hypothesis or number of the counted statement or of its proof was altered. Layout and class: the article's own class is \texttt{amsart} with packages including tikz-cd, amsmath, fourier, amssymb, amscd, amsthm, amsfonts, newlfont, graphicx, mathrsfs, latexsym, tikz, verbatim, xy; the wrapper is the lane's pinned \texttt{amsart} editorial wrapper at 10pt on A4, which restates the theorem-like environments the retained text needs and adds the article's own macro definitions that the retained text needs (none), with extra wrapper packages (bm, bbm, dsfont, xspace, cancel, mathtools). The delivered numbering is the unit's own. Assets: no retained span contains a figure environment, a graphics-inclusion command, a PDF-page inclusion, a table or a TikZ picture; no image file is copied into this package, the package declares no asset dependency and no image file is redistributed with it. Licence: version arXiv:2411.12579v2 of this article is distributed under the Creative Commons Attribution 4.0 International licence, as recorded in the arXiv licence metadata for this exact version and shown on the abstract page https://arxiv.org/abs/2411.12579v2. A full rights notice is delivered at licenses/arxiv-2411.12579-CC-BY-4.0.txt. Characters: every retained excerpt is pure ASCII, so the delivered engine, XeLaTeX with its default Latin Modern text font, reports no missing character.
\begin{lemma} \label{lemma: invariance}
For all $ p, q \geq 0 $, $ f \in \mathfrak{P}_{p,q}(\mathbb{S}^{n-1}) $, and $ U \in \mathcal{U}_n $, we have that $ f \circ U \in \mathfrak{P}_{p,q}(\mathbb{S}^{n-1}) $. The same result holds if\,
$\mathfrak{P}_{p,q}(\mathbb{S}^{n-1}) $  is replaced by $ \mathcal{H}_{p,q}(\mathbb{S}^{n-1}) $.
\end{lemma}

\begin{proposition} \label{propleg-comp}
The polynomial  $L_{n,p,q} \colon \mathbb{C}^n \to \mathbb{C} $ is the unique  polynomial
$g \in \mathcal{H}_{p,q}(\mathbb{C}^n) $ which fulfills the following two properties\,:
\begin{itemize}
\item[(a)]
$g\circ A = g $ for all  $A\in \mathcal{U}_n(e_1)$\,,
\item[(b)]
$g(e_1) = 1 $\,.
\end{itemize}
\end{proposition}

\begin{equation} \label{repro}
  L_{n,p,q}^{\diamond}(w) :=
    \sum_{j=0}^{p\wedge q} c_j(n,p,q) \,w^{p-j} \overline{w}^{q-j}
    (1 - |w|^2)^{j}.
\end{equation}

\begin{equation}\label{integration}
\int_{\mathbb{S}^{n-1}} f(\langle \eta , e_1 \rangle) \, d\sigma_n(\eta) =
\frac{n-1}{\pi} \int_{0}^{1} \int_{-\pi}^{\pi} (1-r^2)^{n-2} f(r e^{i\theta}) \, r \, d\theta \,dr \,.
\end{equation}

\begin{theorem} \label{Koornwinder}
  Let  $p,q \geq 0 $ and  $n \geq 2 $. Then for all  $z = r e^{i \theta} $ with  $r\geq0 $,
  \[
  L^{\diamond}_{n,p,q}(z) = r^{|p-q|}  e^{i(p-q)\theta}
  \,\,
  \frac{(p\wedge q)!(n-2)!}{((p\wedge q) + (n-2))!}
  \,\,
  P_{p\wedge q}^{n-2,|p-q|}(2r^2 -1)\,.
  \]
\end{theorem}

\begin{proof}
Fixing  $n \geq 2$, we   in the following  divide the proof  into two steps.  In the first step we prove that for all  $z = re^{i\theta} \in \overline{\mathbb{D}} $  and  $p,q \ge 0$
  \[
  L^{\diamond}_{n,p,q}(z) = r^{|p-q|}  e^{i(p-q)\theta}
  \,\,
  Q_{p,q}(2r^2 -1) \,,
  \]
 where   $Q_{p,q}: \mathbb{R} \to \mathbb{R} $  is a polynomial  of degree
  less than or equal to $ p \wedge q  $.
  To see this, assume  without loss of generality  that  $p \wedge q = p $.
Then by Equation~\eqref{repro}
  \begin{equation*}
  L_{n,p,q}^{\diamond} (z) =  r^{|p-q|}  e^{i (p-q)\theta}
    \sum_{j=0}^{p}c_j \,r^{2(p-j)}
    (1-r^2)^j  \,,
 \end{equation*}
 and for each  $0 \leq j \leq p $
 \begin{align*}
   &
   r^{2(p-j)} = 2^{-(p-j)}((2r^2 -1)+1)^{p-j} = u_j(2r^2 -1)
   \\&
   (1-r^2)^j
    = (-1)^j 2^{-j}(2(r^2-1))^j = (-1)^j 2^{-j}((2r^2-1)-1))^j =v_j(2r^2 -1)\,,
     \end{align*}
 where  $u_j(x) = 2^{-(p-j)}(x+1)^{p-j} $ is a polynomial of degree  $ p-j $
 and  $v_j(x) = (-1)^j 2^{-j}(x-1)^{j} $ a~polynomial of degree  $ j $.  This clearly proves the claim.




  In the second step we  check that   for all  $r \in [0,1] $ and $p,q \ge 0 $
  \[
  Q_{p,q}(2r^2 -1) = \frac{1}{P_{p\wedge q}^{n-2,|p-q|}(1)}\,\, P_{p\wedge q}^{n-2,|p-q|}(2r^2 -1)\,.
  \]
  To do so, we fix some  arbitrary $k \in \mathbb{Z} $,  and
  show that the preceding  equality holds
  for all  $p,q \ge 0 $ with $k = p-q$.  The idea is to
  prove that in this case all  polynomials $Q_{p,q}$ are
  orthogonal with respect to the weight function
 $ (1-\pmb{\cdot})^{n-2}(1+ \pmb{\cdot})^{|k|}$  on  $[-1,1]$ , which then implies the conclusion using the uniqueness of the Jacobi polynomials.

 More precisely,  assume that for each  $p,q,p',q' \ge 0$
 such that $k = p-q=p'-q'$ and
 $\ell = p\wedge q \neq \ell' = p'\wedge q' $,
 we have,
  \begin{align} \label{koorwin}
    \int_{-1}^{1}  (1-s)^{n-2}(1+s)^{|k|} Q_{p,q}(2s^2 -1)  Q_{p',q'}(2s^2 -1) ds = 0\,.
  \end{align}
Then, by the uniqueness properties of Jacobi polynomials, there exists a constant $c \in \mathbb{R}$ such that
\[
Q_{p,q}(2s^2 -1) = c \, P_{p \wedge q}^{n-2, |k|}(s), \quad s \in [-1,1].
\]
Moreover, since $Q_{p,q}(1) = L_{n,p,q}^{\diamond}(1) = 1$, it follows that
\[
c = \frac{1}{P_{p \wedge q}^{n-2, |p-q|}(1)}.
\]
  This would complete the proof of our claim.

It remains to prove Equation~\eqref{koorwin}. We use Equation~\eqref{integration} to see that
  \begin{align*}
    &
  \int_{\mathbb{S}^{n-1}}
    L_{n,p,q}^{\diamond} (\langle e_1, \eta \rangle)  \overline{
L_{n,p',q'}^{\diamond}(\langle e_1, \eta \rangle)}
    \,d\sigma_n(\eta)
    \\&
  =
  \frac{n-1}{\pi}\int_{0}^{1} \int_{-\pi}^{\pi}
    r^{|p-q|+|p'-q'|}  e^{i (p-q)\theta} e^{-i (p'-q')\theta} (1-r^2)^{n-2} r
    Q_{p,q} (2r^2 -1)  Q_{p',q'} (2r^2 -1)
  d\theta dr
  \\&
  =
  (n-1)\int_{0}^{1}
    (r^2)^{|k|}  (1-r^2)^{n-2} r
    Q_{p,q} (2r^2 -1)  Q_{p',q'} (2r^2 -1)
  dr
  \\&
  =
  2^{-|k|-(n-2)} (n-1)  \int_{-1}^{1} (1+s)^{|k|} (1-s)^{n-2} Q_{p,q}(s) Q_{p',q'} (s) ds\,.
    \end{align*}
where we have used the change of variables  $s=2r^2-1 $ to get the last equality.
    But
      \begin{align*}
      \int_{\mathbb{S}^{n-1}}
    L_{n,p,q}^{\diamond} (\langle e_1, \eta \rangle)  \overline{
L_{n,p',q'}^{\diamond} (\langle e_1, \eta \rangle)}
    \,d\sigma_n(\eta) = 0\,,
    \end{align*}
  since by Proposition~\ref{propleg-comp} we have that
   $L_{n,u,v}^{\diamond} (\langle e_1, \cdot \rangle) \in \mathcal{H}_{u,v}(\mathbb{S}^{n-1})  $ for all  $u,v \geq 0 $, and by Lemma~\ref{lemma: invariance}  that all spaces
   $\mathcal{H}_{u,v}(\mathbb{S}^{n-1})  $ are pairwise orthogonal. This  completes the argument.    \end{proof}

\end{document}
